\chapter{Introducción: qué hace un trader profesional}\label{cap:intro}

Un \emph{trader} profesional no es alguien que adivina hacia dónde va el precio. Es alguien que
\textbf{toma riesgo de forma deliberada, medida y repetible} a cambio de una expectativa positiva, y que
sobrevive lo suficiente como para que esa expectativa se materialice. La diferencia con el aficionado
casi nunca está en la idea de entrada; está en el tamaño, en la disciplina para cortar pérdidas, en los
costes y en el registro honesto de lo que funciona.

\section{Los tipos de trader}
\begin{description}
\item[Sell-side (bancos, creadores de mercado).] Dan liquidez a clientes y ganan el diferencial
  compra-venta (\emph{bid-ask}). Su problema principal es gestionar el inventario y cubrir el riesgo
  que acumulan; sus límites se fijan en DV01, delta, vega y VaR.
\item[Buy-side (fondos, gestoras, hedge funds).] Toman posiciones direccionales o de valor relativo para
  generar rentabilidad. Se miden por Sharpe, drawdown y correlación con el mercado.
\item[Prop trading y traders independientes.] Arriesgan capital propio (o de una firma que les fondea)
  con reglas de pérdida diaria y \emph{drawdown} máximo muy estrictas.
\item[Sistemáticos (\emph{quants}).] Codifican las reglas y dejan que un ordenador las ejecute; su
  ventaja está en la investigación y en la ejecución, y su riesgo en el sobreajuste.
\end{description}

\section{La ecuación del trader}
Si cada operación arriesga una unidad $1R$ (la pérdida si salta el stop), la \emph{esperanza} por
operación es
\begin{equation}\label{eq:esperanza}
  E = p\cdot \bar G - (1-p)\cdot \bar P \qquad (\text{en } R),
\end{equation}
con $p$ la tasa de acierto, $\bar G$ la ganancia media y $\bar P$ la pérdida media (en $R$). El
resultado anual es aproximadamente $E \times N \times (\text{riesgo por operación})$, con $N$ el número
de operaciones. Todo el curso trata de mejorar uno de esos tres factores sin destruir los otros:
\begin{itemize}
\item subir $E$: mejores ideas, mejor ejecución, menos costes (capítulos~\ref{cap:micro}, \ref{cap:tecnico}, \ref{cap:equity}--\ref{cap:cripto});
\item que $E$ sea \emph{real} y no un artefacto del backtest (capítulo~\ref{cap:backtest});
\item elegir el riesgo por operación que maximiza el crecimiento sin arruinarse (capítulo~\ref{cap:riesgo});
\item ejecutar el plan sin que las emociones lo rompan (capítulo~\ref{cap:psico}).
\end{itemize}

\begin{clave}
Una tasa de acierto del 35~\% puede ser excelente si las ganancias medias son 3 veces las pérdidas
($E = 0{,}35\cdot3 - 0{,}65 = 0{,}40R$). Un 80~\% de acierto puede ser ruinoso si cada pérdida es 5
veces la ganancia media. \textbf{El acierto, solo, no significa nada.}
\end{clave}

\section{Cómo usar este curso}
Cada capítulo combina teoría, fórmulas y un enlace a funciones de la librería \py{trader}
(carpeta \py{curso\_trader/}). Las fórmulas que ves aquí son las que están implementadas y
verificadas por los tests. Para instalar y comprobar:
\begin{lstlisting}
cd curso_trader
pip install -r requirements.txt
pytest tests                      # ~15 s, sin red
python scripts/generar_figuras.py # regenera las figuras de este PDF
\end{lstlisting}
Los scripts de \py{scripts/} están pensados para el día a día (capítulo~\ref{cap:rutina}):
tamaño de posición, hoja premercado, informe de riesgo, backtest, diario de operaciones y
calculadoras por clase de activo. Todos funcionan en modo demostración con datos sintéticos y aceptan
tus propios CSV.

\paragraph{Estructura.} La parte~I cubre lo común a todos los mercados (microestructura, estadística,
análisis técnico, riesgo, backtesting y psicología). La parte~II recorre las clases de activo:
\textbf{renta variable (equity)}, \textbf{divisas (FX)}, \textbf{materias primas (commodities)}, y además
\textbf{renta fija}, \textbf{opciones y volatilidad} y \textbf{cripto}, porque un trader de cualquier
mesa necesita entender tipos de interés y opciones. La parte~III convierte todo en una rutina diaria y
un plan de estudio de 12 semanas.
